%% ---------------------------------------------------------------------
\section{El modelo de Kuramoto}
\label{sec:oc-kuramoto}
%% ---------------------------------------------------------------------

\notamargen{Esta sección es fenomenológica: se describe lo que pasa y por qué es plausible. Las pruebas de sincronización usan herramientas de Lyapunov y de pasividad (capítulos~7 a~11).}
Kuramoto tomó \eqref{eq:oc-redsin} con todos contra todos y el mismo peso:
\begin{equation}
  \dot\theta_i=\omega_i+\frac KN\sum_{j=1}^N\sin(\theta_j-\theta_i),\qquad i=1,\dots,N,
  \label{eq:oc-kuramoto}
\end{equation}
donde $K\ge0$ es la intensidad del acoplamiento y las frecuencias naturales $\omega_i$ se sacan de una distribución con densidad $g(\omega)$, simétrica alrededor de su centro (que, cambiando de marco de referencia, se puede tomar como cero). La pregunta: ¿a partir de qué $K$ emerge la sincronía?

Para medirla se usa el \clave{parámetro de orden}
\begin{equation}
  r\,e^{j\Psi}=\frac1N\sum_{j=1}^Ne^{j\theta_j} .
  \label{eq:oc-orden}
\end{equation}
Es el centroide de las $N$ fases vistas como puntos del círculo unitario: $r\in[0,1]$ mide la coherencia ($r\approx0$, fases repartidas; $r=1$, todas iguales) y $\Psi$ es la fase promedio. Multiplicando \eqref{eq:oc-orden} por $e^{-j\theta_i}$ y tomando la parte imaginaria, $r\sin(\Psi-\theta_i)=\frac1N\sum_j\sin(\theta_j-\theta_i)$, y el modelo queda
\begin{equation}
  \dot\theta_i=\omega_i+K\,r\sin(\Psi-\theta_i) .
  \label{eq:oc-campomedio}
\end{equation}
Cada oscilador ya no ve a los demás uno por uno: ve un solo ``estímulo'' de fase $\Psi$ e intensidad $Kr$. Es la luciérnaga frente a la linterna, \eqref{eq:oc-dos}, con la diferencia de que la linterna la fabrican entre todos. El mecanismo es una retroalimentación: un poco de coherencia ($r>0$) hace que más osciladores se enganchen (los que cumplen $\abs{\omega_i}\le Kr$), y eso aumenta $r$.

Kuramoto mostró que, cuando $N\to\infty$, la solución incoherente $r=0$ pierde la estabilidad en
\begin{equation}
  K_c=\frac2{\pi g(0)},
  \label{eq:oc-Kc}
\end{equation}
y que, para la distribución de Lorentz $g(\omega)=\frac{\gamma}{\pi(\gamma^2+\omega^2)}$, la coherencia de régimen es exactamente $r=\sqrt{1-K_c/K}$ para $K>K_c=2\gamma$ \citep{Kur84}. La \cref{fig:oc-kuramoto} compara esa predicción con una simulación de $N=400$ osciladores con $\gamma=0.5$: $r=0.58$, $0.71$ y $0.82$ para $K=1.5$, $2$ y $3$, contra $0.577$, $0.707$ y $0.816$ de la teoría.

\video{V5_1_kuramoto}{Las fases se agrupan; algunas siguen dando vueltas}
\script{cap05_orbitas_ciclos/C5_Ej6_Kuramoto.m}
\begin{figure}[ht]
\centering
\includegraphics[width=\linewidth]{cap05/f_kuramoto}
\caption{Kuramoto con $N=400$ y frecuencias de Lorentz ($\gamma=0.5$, $K_c=1$). Las fases en $t=100$ (color: frecuencia natural) para $K=0.5$ y $K=2$, con la flecha $re^{j\Psi}$. Derecha: $r$ de régimen contra la predicción $\sqrt{1-K_c/K}$.}
\label{fig:oc-kuramoto}
\end{figure}

\codigorepo{cap05_orbitas_ciclos/C5_Ej6_Kuramoto.m}{kuramoto}{Parámetro de orden de régimen para varios $K$}

\begin{trampa}[label=tr:oc-kuramoto]{Acoplar no es igualar}
\plausible{Con $K$ suficientemente grande, todas las fases convergen a la misma: $r\to1$.}
Con frecuencias naturales distintas, las fases \emph{no} pueden ser iguales: si $\theta_i=\theta_j$ para todo $i,j$, entonces $\dot\theta_i=\omega_i$ y se separan de inmediato. Lo que se alcanza es un estado en que los osciladores enganchados giran a una frecuencia común con desfases fijos, y para Lorentz quedan además osciladores rápidos sin enganchar (los de $\abs{\omega_i}>Kr$), que siguen dando vueltas. Aun con $K=3$, $r=0.82$.
\correcto{Para $K>K_c$, $r$ crece con $K$ pero se queda por debajo de $1$; solo $r\to1$ cuando $K\to\infty$. Igual que en la \cref{tr:oc-enfase}: el acoplamiento da sincronía de frecuencia, no de fase, salvo que las frecuencias sean idénticas.}
\end{trampa}

\begin{pausa}
Con dos osciladores ($N=2$) y frecuencias $\pm\Delta/2$, ¿qué valor de $K$ en \eqref{eq:oc-kuramoto} los engancha?
\respuesta{Con $N=2$, $a_{12}=a_{21}=K/2$ en \eqref{eq:oc-redsin}, y \eqref{eq:oc-dos} da $\dot\varphi=\Delta-K\sin\varphi$. Hay enganche si y solo si $K\ge\Delta$.}
\end{pausa}

\begin{profundiza}[Investigación: sincronización resiliente con pocos actuadores]
En una red real (una red eléctrica, una red de sensores) hay perturbaciones y se puede actuar sobre algunos nodos:
\[
  \dot\theta_i=\omega_i+\frac KN\sum_{j=1}^Na_{ij}\sin(\theta_j-\theta_i)+u_i(t)+d_i(t),
\]
con $d_i$ una disrupción y $u_i$ un control. ¿Cómo recuperar la sincronía después de una disrupción controlando solo unos pocos nodos y gastando poca energía? Es la pregunta del proyecto de Chacón-Chamorro, Charria-Macías, Giraldo y Quijano sobre sincronización resiliente de redes de Kuramoto \citep{CCGQ26}: una arquitectura en la que una red neuronal escoge en línea qué nodos actuar (\emph{pinning}) y un controlador de consenso los lleva de vuelta al estado sincronizado.
\end{profundiza}
